\chapter{Training Stability}
\label{chap:stability}
Now that we have covered architectures and optimisation methods, in theory we can design models which converge.
However, in practice, these things are very delicate and we need to carefully tune our design according to our chosen architecture, optimisation method and other design choices.

\section{Initialisation}
\label{sec:Initialisation}
A deep network repeatedly applies an update, and so an intuitive question is \textit{where should we start?}.
Because forward activations and backward updates can grow or shrink layer after layer, any instabilities compound and explode/zero, rather than cancelling out.

Zero-initialising all weights in a symmetric hidden layer prevents neurons from differentiating because equivalent units receive equivalent updates.
So, we introduce randomness to break the symmetry, and deploy variance scaling to keep the starting signal numerically usable.

\subsection{Xavier / Glorot and He / Kaiming Initialisation}
\label{sec:stability-xavier-glorot-and-he-kaiming-initialisation}
For a dense layer, common Gaussian starting variances are
\begin{align}
    \theta_{ij,\mathrm{Xavier}}
    &\sim
    \mathcal N\!\left(0,
    \underbrace{\frac{2}{N_{\text{inputs}}+N_{\text{outputs}}}}_{\text{balance forward/backward scale}}
    \right),\\
    \theta_{ij,\mathrm{He}}
    &\sim
    \mathcal N\!\left(0,
    \underbrace{\frac{2}{N_{\text{inputs}}}}_{\text{compensate ReLU half-zeroing}}
    \right).
\end{align}
Xavier is a common starting point for linear/tanh-like activations; He/Kaiming is designed around ReLU-like activations.
These are approximate signal-scale arguments, not guarantees that training remains well conditioned.

\begin{calloutpractical}
    ReLU-like activation: start with He/Kaiming; linear or tanh-like activation: Xavier/Glorot is a common baseline.
\end{calloutpractical}

\section{Normalisation}
\label{sec:normalisation}
Normalisation methods mainly differ in \textbf{which axes define the statistic}.
\figpdf[.89\textwidth]{normalization_axes}{BatchNorm and LayerNorm normalise different slices of an activation tensor.}

For $N_{\text{values}}$ values in the selected normalisation group:
\[
    \underbrace{\hat u_i}_{\text{centred/scaled value}}
    =
    \frac{
     u_i-\overbrace{\frac{1}{N_{\text{values}}}\sum_j u_j}^{\mu_{\text{group}}}
    }{
    \sqrt{
    \underbrace{\frac{1}{N_{\text{values}}}\sum_j(u_j-\mu_{\text{group}})^2}_{\sigma^2_{\text{group}}}
    +\epsilon
    }}.
\]
A learned scale and optional shift can restore useful amplitude/offset after normalisation.

\subsection{Batch Normalisation}
\label{subsec:stability-batch-normalisation}
BatchNorm averages across examples for corresponding features/channels during training and usually uses running training statistics at inference.
It can improve optimisation, but tiny or nonrepresentative batches make its reference statistics noisy or mismatched.
For example: In a house price prediction model, we normalise the \texttt{neighbour\_house\_price} feature to the average of the other houses in the batch.

\subsection{Layer Normalisation}
\label{subsec:stability-layer-normalisation}
LayerNorm averages within each example/token across a chosen feature set, so its statistics do not depend on unrelated examples in the batch.
This is why it fits sequence models with variable or small batch sizes well.
For example, In a house price prediction model, we may average all the input features of one example.

\subsection{RMSNorm}
\label{subsec:stability-rmsnorm}
RMSNorm controls vector magnitude without subtracting the mean:
\[
    \operatorname{RMSNorm}(h)_i
    =
    \theta_{\text{scale},i}
    \frac{h_i}{
    \sqrt{\frac{1}{N_{\text{features}}}\sum_j h_j^2+\epsilon}}
    .
\]
It is cheaper than full centring-and-scaling and is common in modern language models.

\begin{table}[htbp]
    \centering
    \small
    \begin{tabularx}{\linewidth}{@{}L{.19\linewidth}YYY@{}}
    \toprule
     & BatchNorm & LayerNorm & RMSNorm \\
    \midrule
    Statistics over & Batch, per feature/channel & Features within one example/token & Features within one example/token \\
    Subtract mean? & Yes & Yes & No \\
    Scale by & Standard deviation & Standard deviation & Root mean square \\
    Depends on other examples? & During training, yes & No & No \\
    \bottomrule
    \end{tabularx}
    \caption{Normalisation methods.}
\end{table}

These methods do not generally whiten cross-feature covariance; whitening is separated in \cref{sec:representation-whitening-pca-and-zca,sec:representation-whitening-versus-normalisation}.

\section{Regularisation}
Another method of improving stability is to penalise unstable changes within the model. 
For example, we may want to penalise parameter updates which are very large, because these \textit{can} lead to gradient explosions.  
\label{sec:stability-regularisation-l1-and-l2}
A regularised objective exposes the trade-off directly:
\[
    \underbrace{\mathcal L_{\text{total}}(\theta)}_{\text{what optimisation sees}}
    =
    \underbrace{\mathcal L_{\text{data}}(\theta)}_{\text{fit observations}}
    +
    \underbrace{\lambda R(\theta)}_{\text{parameter preference}}.
\]
L2 uses $R(\theta)=\|\theta\|_2^2$ and usually shrinks parameters smoothly.
L1 uses $R(\theta)=\|\theta\|_1$ and can drive coefficients exactly to zero.
Elastic Net combines both.
The interpretation depends on parameterisation, so ``small weights'' is not universally the same as ``simple realised function''.

\begin{calloutintuition}
    L1 is basically feature selection when individual weights map cleanly to input features.
    L2 shrinks rather than deletes them.
\end{calloutintuition}

The Bayesian MAP view is in \cref{sec:foundations-maximum-likelihood-and-maximum-a-posteriori-estimation}.
For Adam, decoupled weight decay changes the update relative to putting an L2 gradient inside Adam; see \cref{sec:optimisation-adamw-decouple-weight-decay-from-gradient-adaptation}.

\subsection{Dropout}
\label{sec:stability-dropout-and-stochastic-regularisation}
In inverted dropout, each activation is randomly kept with probability $1-p_{\text{drop}}$ and rescaled so its conditional expectation is unchanged:
\[
    \tilde h_j
    =
    \underbrace{\frac{m_j}{1-p_{\text{drop}}}}_{\text{random keep + expectation-preserving scale}}h_j,
    \qquad
    m_j\sim\operatorname{Bernoulli}(1-p_{\text{drop}}).
\]
Dropout discourages dependence on one exact active pathway, but too much stochasticity can make optimisation unnecessarily difficult.
It is disabled at inference.

\section{Augmentation}
\label{sec:stability-data-augmentation-and-symmetry-assumptions}
Augmentation asserts that a transformation should preserve or predictably transform the target.
A horizontal flip may preserve animal identity but invert left/right anatomy, time reversal may violate causality in a physical signal.
Augmentation is an inductive-bias choice, not free extra data.
See \cref{sec:nfl-inductive-bias-the-assumptions-that-make-learning-possible}.



